Электрические цепи постоянного тока являются фундаментальной основой современной электротехники, электроники и энергетики. Практически любое электротехническое устройство — от простейшего бытового прибора до сложной промышленной установки — может быть представлено в виде эквивалентной электрической схемы, состоящей из источников питания и резистивных элементов. Умение выполнять расчёт таких цепей, определять токи, напряжения и мощности в каждой ветви, а также проверять баланс мощностей представляет собой одну из ключевых компетенций инженера-электрика.

Особый интерес представляют разветвлённые электрические цепи, содержащие как последовательные, так и параллельные участки соединения резисторов. Расчёт подобных цепей требует применения основных законов электротехники: закона Ома, первого и второго законов Кирхгофа, а также метода эквивалентных преобразований. Важным этапом любого расчёта является проверка баланса мощностей, которая позволяет оценить корректность выполненных вычислений и убедиться в соблюдении закона сохранения энергии.

Целью настоящей работы является разработка программного приложения на языке C\# для расчёта параметров разветвлённой электрической цепи постоянного тока со смешанным соединением резисторов (последовательно-параллельная схема).

Для достижения поставленной цели необходимо решить следующие задачи:
\begin{enumerate}
    \item Провести анализ электрической схемы и определить математическую модель цепи.
    \item Выбрать и обосновать методы расчёта (закон Ома, правила Кирхгофа, эквивалентные преобразования). 
    анализаторов.
    \item Реализовать на языке C\# алгоритм расчёта токов, напряжений и мощностей для всех элементов цепи.
    \item Реализовать проверку баланса мощностей и законов Кирхгофа.
\end{enumerate}